\section{Hashing con orden}

\begin{frame}
  \vspace{1.5cm}
  \centering
  {\LARGE Hashing con orden}
\end{frame}

\begin{frame}{Comprimir}
  Hay que hacer \(N^{2}\) comparaciones entre \(N\) strings, con \(N = 2000\) y cada string de largo \(2000\).

  \medskip
  Se ordenan una vez. Cada string pasa a ser su índice en ese orden, y comparar es comparar enteros.

  \vspace{0.7em}
  \begin{center}
  \begin{tikzpicture}[font=\small]
    \foreach \n/\x/\word/\id in {a/0/zorro/2, b/2.3/gato/0, c/4.6/perro/1, d/6.9/gato/0} {
      \node (\n) at (\x,0) {\texttt{\word}};
      \node[circle, draw=shared, thick, fill=shared!10, minimum size=7mm, inner sep=0pt] (i\n) at (\x,-1.35) {\id};
      \draw[e] (\n) -- (i\n);
    }
  \end{tikzpicture}
  \end{center}

  \vspace{0.2em}
  Orden de los distintos: \texttt{gato}, \texttt{perro}, \texttt{zorro}.
  Eso es comprimir: ordenar, borrar repetidos, reemplazar por el índice.
\end{frame}

\begin{frame}{Concatenar ids}
  \(A\) son strings de largo \(3\), \(B\) de largo \(2\).
  Los ids conservan el orden adentro de cada conjunto.

  \vspace{0.45em}
  \begin{columns}[c,onlytextwidth]
    \begin{column}{0.42\textwidth}
      \centering
      \begin{tabular}{@{}l@{\ \ }c@{\ \ }c@{}}
        \(A\) & \texttt{gat} & \texttt{perr} \\
            & \(0\) & \(1\) \\[0.45em]
        \(B\) & \texttt{as} & \texttt{os} \\
            & \(0\) & \(1\)
      \end{tabular}
    \end{column}
    \begin{column}{0.54\textwidth}
      \centering
      \renewcommand{\arraystretch}{1.15}
      \begin{tabular}{@{}c@{\quad}l@{\quad}c@{}}
        par & concatenación & id \\
        \((0,0)\) & \texttt{gatas} & \(0\) \\
        \((0,1)\) & \texttt{gatos} & \(1\) \\
        \((1,0)\) & \texttt{perras} & \(2\) \\
        \((1,1)\) & \texttt{perros} & \(3\)
      \end{tabular}
    \end{column}
  \end{columns}

  \vspace{0.55em}
  El orden de los pares es el orden de las concatenaciones.
  Esos pares se comprimen, y los ids nuevos quedan ordenados como las strings \(ab\).
\end{frame}

\begin{frame}{Concatenar ids}
  El corte entre las dos partes tiene que caer siempre en el mismo lugar.
  Si los largos de \(A\) no son todos iguales, o los de \(B\), el par deja de representar a la concatenación.

  \vspace{0.6em}
  \begin{center}
  \renewcommand{\arraystretch}{1.35}
  \begin{tabular}{@{}l@{\quad}c@{\quad}c@{\quad}l@{}}
    \texttt{a} \(\to 0\), \texttt{ab} \(\to 1\) &
      \texttt{a}+\texttt{bc} & \texttt{abc} & par \((0,0)\) \\
    \texttt{bc} \(\to 0\), \texttt{c} \(\to 1\) &
      \texttt{ab}+\texttt{c} & \texttt{abc} & par \((1,1)\)
  \end{tabular}
  \end{center}

  \vspace{0.45em}
  Las dos concatenaciones son la misma string.
  Los pares no, y al comprimirles salen ids distintos.
\end{frame}
